\section{Every run as a Job and as a Flyte task}
\label{sec:twins}

The workload code was generated as Flyte tasks, but nothing in the body of a task
depends on Flyte. We use that to put every measured run into the benchmark twice, once as
a raw Kubernetes Job and once as a Flyte task, with the task body unchanged in both. To
make the Job version, \code{job.py}, we turn \code{File} parameters into paths, drop the
\code{flyte} import, and add a \code{\_\_main\_\_} block that stages the inputs, calls the
function and writes \code{result.json} for upload. A \code{batch/v1} manifest goes with
it, carrying the image, the input locations, the thread settings the measurement used,
the staging volume and an empty \code{resources} block, which is the request still to be
decided. \Cref{fig:twin} shows one task written both ways.

\begin{figure}[tbp]
\centering
\begin{minipage}[t]{0.48\linewidth}
\textbf{\small Flyte task}
\lstinputlisting[language=Python]{figures/ex_flyte.py}
\textbf{\small Job manifest (boto3-tmpfs)}
\lstinputlisting{figures/ex_manifest.yaml}
\end{minipage}\hfill
\begin{minipage}[t]{0.50\linewidth}
\textbf{\small \code{job.py} (boto3 staging)}
\lstinputlisting[language=Python]{figures/ex_job.py}
\end{minipage}
\caption{One RightSizeMem task, as a Flyte task and as a raw Job. The disk staging
differs from this manifest in one line, \code{emptyDir: \{\}}. The init-container
stagings add an \code{aws-cli} init container that copies the inputs to \code{/data}, and
\code{job.py} then reads them from \code{argv} (\cref{app:prompt}).}
\label{fig:twin}
\end{figure}

\subsection{Stagings}
How a Job's inputs reach its container is the biggest thing that changes its memory
need. Each task is assigned one of four stagings by a hash of its id, and all of its runs
share it (\cref{tab:stagings}). Inputs are fetched either by an init container running
\code{aws-cli} or by the job itself with boto3, and they land either on node disk or on a
memory-backed volume.

\begin{table}[tbp]
\centering\small
\caption{The four Job stagings.}
\label{tab:stagings}
\resizebox{\linewidth}{!}{%
\begin{tabular}{@{}lrlll@{}}
\toprule
staging & share & who fetches the inputs & where they land & memory need beyond the workload \\
\midrule
init-disk   & 30\% & aws-cli init container & \code{emptyDir} on node disk & --- \\
init-tmpfs  & 20\% & aws-cli init container & \code{emptyDir}, \code{medium: Memory} & $+B$ (staged bytes) \\
boto3-disk  & 30\% & the job itself, boto3  & \code{emptyDir} on node disk & $+48$\,MiB (boto3) \\
boto3-tmpfs & 20\% & the job itself, boto3  & \code{emptyDir}, \code{medium: Memory} & $+48$\,MiB $+B$ \\
\bottomrule
\end{tabular}}
\end{table}

\subsection{Label accounting}
\label{sec:accounting}
Let $w$ be the workload's measured peak memory (\cref{sec:dataset}) and let $B$ be the
size of the staged inputs plus the result file, with each file rounded up to whole pages.
The two kinds of pod need
\begin{align*}
\text{Job:}\quad & 1.05\,\bigl(w + 48\,\text{MiB}\cdot[\text{boto3}] + B\cdot[\text{tmpfs}]\bigr),\\
\text{Flyte:}\quad & 1.05\,\bigl(w + R\bigr), \qquad R = 70\,\text{MiB},
\end{align*}
where the factor 1.05 is headroom and the memory label is the smallest grid cell at or
above the need. Each term has a simple reason. Files on a memory-backed volume are
charged to memory byte for byte, while the same files on node disk are reclaimable page
cache and add nothing, since the measured workloads already read local files. The boto3
client adds 48\,MiB to a workload's peak. Its download also runs inside the job, taking
$d = \text{bytes}/440\,\text{MB/s} + 0.4\,\text{s per object} + 0.7\,\text{s}$ at about
1.8 cores, so we prepend that window to the run's CPU series. For a light job with a large
input, the download can be the busiest 15 seconds of the run. The Flyte runtime $R$ is
70\,MiB, measured in real task pods (\cref{sec:validation}). Earlier releases used a value
of 180\,MiB that had never been measured. Staging does not change GPU need.

The manifest sizes the init container at $256\,\text{MiB} + B$ for tmpfs stagings, so
that it can hold what it writes. Our first release used a flat 256\,Mi, which real pods
showed cannot work for larger inputs (\cref{sec:validation}). Histories are reported in
each interface's own units: a Job's earlier runs report the container's peak as cAdvisor
would, and a Flyte task's report the peak including the runtime, as Flyte's UI would.

The staging constants come from an in-cluster probe using boto3 1.43 under cgroup~v2. A
bare interpreter holds 11\,MiB. With boto3 imported, a client built and a file downloaded,
the process holds 74--80\,MiB, and it never peaks above that final level. Inside a real
workload boto3 adds only 48\,MiB at the median, because it shares modules the workload has
already loaded. Download throughput rises from 122\,MB/s for a 64\,MiB object to
417--431\,MB/s for a 4\,GiB one.
